\documentclass[11pt,a4paper]{article}
\usepackage{../shared/luxcover}
\usepackage[utf8]{inputenc}
\usepackage[T1]{fontenc}
\usepackage{amsmath,amssymb}
\usepackage{hyperref}
\usepackage{booktabs}
\usepackage{enumitem}
\usepackage{xcolor}
\usepackage{tabularx}

\hypersetup{colorlinks=true,linkcolor=black,citecolor=blue,urlcolor=blue}

\title{Lux Industries, Inc.\\[4pt]
\large Investor Memorandum: \$100M Growth Round\\
Post-Quantum Sovereign Blockchain, the W3A IP Portfolio,\\
and the Regulated-Exchange Launch}

\author{
Zach Kelling\thanks{Corresponding author: zach@lux.network}\\
\textit{Lux Industries, Inc.}\\
\texttt{research@lux.network}
}

\date{July 2026 \quad|\quad CONFIDENTIAL --- PREPARED FOR ZACH FRIEDMAN}

\begin{document}

\luxcoverpage

\begin{abstract}
\noindent This memorandum invites Zach Friedman to lead a \$100~million
growth round in Lux Industries, Inc. Mr.\ Friedman is uniquely positioned to
lead: he is a co-founder of \textbf{GDA Capital} (a blockchain merchant bank
and venture studio) and of \textbf{Secure Digital Markets (SDM)}, a global
digital-asset brokerage through which he has facilitated over \$2~billion in
transactions --- and he has already invested personally into a Lux Network
validator, taking protocol-level exposure ahead of this round. This round
converts that alignment into a lead position at the moment three things come
together: a \emph{live} post-quantum public mainnet with on-chain matching
proven in production; a granted, 2018-priority intellectual-property
portfolio being consolidated under the W3A Foundation; and a regulated
white-label exchange stack (ATS / broker-dealer / transfer agent) that turns
the protocol into a licensable revenue engine --- for which SDM is the
natural anchor venue. The use of proceeds is concentrated on the three
value-inflection events that each de-risk the next: closing and
consolidating the IP, taking Lux public with deep liquidity and banking
rails, and standing up the first regulated venues on the sovereign-exchange
platform. Figures herein are a working frame for discussion, not an offer;
terms are subject to definitive documentation and diligence.
\end{abstract}

\tableofcontents
\newpage

%% ============================================================
\section{The Ask}

\begin{itemize}[leftmargin=1.4em]
  \item \textbf{Raise.} \$100~million growth round for Lux Industries,~Inc.
  \item \textbf{Lead.} Zach Friedman invited to lead. His existing personal
  investment into a Lux Network validator is the anchor commitment and the
  conviction signal to the syndicate; the lead allocation and board/observer
  rights are open for negotiation.
  \item \textbf{Instrument \& valuation.} To be set in definitive documents.
  A priced equity round with a token-warrant (or a convertible with a
  network-token side letter) fits a company whose value accrues in both the
  operating entity and the protocol. No valuation is asserted in this memo;
  it is a diligence-and-negotiation output.
  \item \textbf{Use of proceeds.} Three inflection events --- IP
  consolidation (W3A), public launch \& banking, and the regulated-exchange
  rollout --- detailed in \S\ref{sec:use}.
  \item \textbf{Why now.} The mainnet is live and finalizing; the exchange
  stack runs today; the anchor IP asset carries a first-post-closing
  maintenance deadline (Oct 2026, \S\ref{sec:ip}). The window to enter
  before public launch is finite.
\end{itemize}

%% ============================================================
\section{Why Lux, Why Now}

Lux is a post-quantum, multi-chain blockchain platform. The thesis for this
round is not a whitepaper --- it is that the hard, un-fakeable milestones
are already behind us, and capital now buys \emph{scale and distribution},
not \emph{science risk}.

\begin{enumerate}[leftmargin=1.6em]
  \item \textbf{The mainnet is live.} Lux runs a public primary network with
  a production EVM C-Chain and a family of native chains. This is operating
  infrastructure with real state and real validators --- not a testnet.
  \item \textbf{On-chain matching works in production.} The native
  decentralized-exchange precompile settled a live maker/taker fill on
  mainnet: an order was placed, crossed, and settled atomically inside the
  block, moving both legs of the trade in a single accepted transaction with
  exact value conservation. Cryptographic, in-block settlement is
  demonstrated, not promised.
  \item \textbf{Custody is dealerless.} The signing key for custodied assets
  is generated by a verifiable secret-sharing protocol and is never
  assembled at any single party --- no qualified-custodian single point of
  compromise. This is the property regulated venues and banks cannot get
  from an incumbent custodian.
  \item \textbf{The regulated venue is a product, not a plan.} A white-label
  platform --- Alternative Trading System, broker-dealer, and transfer-agent
  services over the native custody / matching / settlement primitives ---
  runs today and is jurisdiction-configurable (U.S.\ and Canadian regimes).
  This is the licensable, recurring-revenue surface.
  \item \textbf{The IP moat is granted and dated.} A 2018-priority,
  three-patent family (two U.S., one European), plus Lux's own
  post-quantum-cryptography portfolio, are being consolidated under one
  holding entity (\S\ref{sec:ip}).
\end{enumerate}

\noindent Capital is the missing input to convert live technology into
market position: liquidity for the public launch, licenses and banking
partnerships for the regulated venues, and the closing of the IP portfolio.

%% ============================================================
\section{Use of Proceeds}
\label{sec:use}

The \$100M is allocated to the three inflection events, in the order in
which each de-risks the next. The split below is an illustrative working
framework subject to board approval and final round sizing --- not a
covenant.

\begin{table}[h]
\centering
\small
\begin{tabularx}{\textwidth}{@{}lrX@{}}
\toprule
\textbf{Bucket} & \textbf{Illustrative} & \textbf{What it funds}\\
\midrule
IP: W3A + portfolio        & \$10--15M & '853 family acquisition, PQ-portfolio consolidation, prosecution + continuations-in-part (\S\ref{sec:ip}), maintenance reserve.\\
Public launch + banking    & \$30--35M & Listings and protocol-owned liquidity, market-making, ecosystem grants, exchange integrations, and the banking / fiat-rail partnerships the venues require.\\
Regulated-exchange rollout & \$30--35M & Licensing (broker-dealer / ATS / TA, per jurisdiction), first white-label venue deployments, compliance and KYC/AML provider integrations, custody operations.\\
Team + working capital     & \$15--20M & Engineering, legal/regulatory, BD, and 18--24 months of operating runway with reserve.\\
\midrule
\textbf{Total}             & \textbf{\$100M} & \\
\bottomrule
\end{tabularx}
\end{table}

The sequencing matters. IP consolidation is fast and cheap relative to the
round and lifts the floor value of everything else. The public launch
creates the liquid asset and the token-side value that the equity is
partly a claim on. The regulated-exchange rollout is where durable,
recurring revenue lives --- each white-label venue is a licensing and
transaction-fee annuity, and the first signed venue is the proof that
underwrites the rest.

%% ============================================================
\section{The IP Portfolio and the W3A Foundation}
\label{sec:ip}

A separate memorandum (\emph{W3A Foundation: Acquisition of U.S.\ Patent
11,803,853 and Consolidation of the Post-Quantum IP Portfolio}) covers this
asset in full; the essentials for this round follow.

\begin{itemize}[leftmargin=1.4em]
  \item \textbf{The anchor family.} ``Controlling transactions on a
  network'' (inventor Mark Briscombe): US~11,803,853~B2, US~12,586,073~B2,
  and EP~3\,788\,757~B1 --- all granted, May~2018 priority, expiring on or
  about 6~May~2039 (\(\approx\)12.8 years of term). Owner of record:
  Nahmii~AS (Bergen, Norway). The family was extended as recently as
  March~2026, when the U.S.\ continuation granted.
  \item \textbf{What it covers.} A network architecture that separates
  transaction origination from execution, with execution gated on
  independent validator + oracle verification --- the control-layer pattern
  every L2 sequencer, cross-chain bridge, and MPC-custody withdrawal flows
  through. The 2018 priority date predates essentially every modern rollup,
  production bridge, and oracle-gated settlement system in use.
  \item \textbf{Two layers, one moat.} Lux's post-quantum work protects the
  \emph{signature} layer (threshold ML-DSA / R-LWE / SLH-DSA, the Quasar
  consensus family, on-chain PQ-verification precompiles). The '853 family
  protects the \emph{control} layer above it. A signature portfolio alone
  invites ``design around the scheme''; a control patent alone invites
  ``design around the oracle.'' Held together under W3A, they cover both
  what signs and what gates.
  \item \textbf{Continuations-in-part.} The strategic upside this round
  funds: filing CIPs that carry the 2018 priority \emph{forward} onto Lux's
  own post-quantum and dealerless-custody implementations --- extending the
  original transaction-control claims onto (i)~post-quantum-signed
  validator/oracle gating, (ii)~threshold-MPC custody withdrawal flows, and
  (iii)~on-chain matching with in-block settlement. This is the mechanism by
  which an acquired 2018 asset becomes a moat around \emph{Lux's} 2026
  production stack, not merely a defensive backstop. (Scope and claimability
  of any CIP are counsel-and-diligence questions; a broadening path also
  remains available on the '073 continuation into 2028.)
  \item \textbf{Structure.} The family is acquired by, and all Lux PQ
  filings assign into, the W3A Foundation --- wholly owned by Lux
  Industries. Operating entities take a perpetual intra-group license back;
  external licensing runs through the foundation. One entity, one licensing
  surface, one place for an investor to value.
  \item \textbf{Timing.} The first U.S.\ maintenance fee on '853 opens
  31~October~2026 --- the first post-closing docket item, and a concrete
  reason the acquisition is time-sensitive.
\end{itemize}

%% ============================================================
\section{The Regulated-Exchange Revenue Engine}

The protocol is the moat; the regulated exchange is the business model.
Lux's white-label platform lets an operator stand up a fully regulated
venue in which custody, matching, settlement, and compliance are native
chain primitives rather than stitched-together vendors:

\begin{itemize}[leftmargin=1.4em]
  \item \textbf{One venue, every asset.} Crypto, tokenized public equities,
  and private offerings (Reg~D / A+ / CF) trade on the same order book,
  under the same dealerless-custody root, recorded in the same on-chain
  register.
  \item \textbf{Jurisdiction as configuration.} The same core instantiates a
  U.S.\ (SEC/FINRA) or a Canadian (CIRO / provincial-commission / FINTRAC)
  venue by swapping a declarative brand profile --- the hard,
  trust-critical layer is built and audited once.
  \item \textbf{Recurring economics.} Each white-label venue is a licensing
  + transaction-fee annuity. This round funds converting the pipeline into
  signed, launched venues.
\end{itemize}

\paragraph{SDM as the anchor venue.}
The natural first deployment is \textbf{Secure Digital Markets} itself. SDM
is a live, institutional digital-asset brokerage that has already
intermediated \$2B+ in transactions --- exactly the operator profile that
the regulated stack is built for. Extending SDM from OTC digital-asset
brokerage into a full tokenized-securities venue (adding ATS matching,
transfer-agent registry, and on-chain dealerless custody) is the highest-fit
first launch in the pipeline, and it is one the lead investor is positioned
to champion from the inside. A flagship venue operated by an aligned
principal turns the ``first customer'' risk --- usually the hardest gap for
a B2B platform --- into a starting condition.

\noindent This is the answer to ``how does the token thesis become a
company'': durable B2B revenue from regulated operators, on top of a
protocol whose usage the equity and token both capture.

%% ============================================================
\section{Zach Friedman: Why This Lead}

A lead investor is worth most when capital is the least of what they bring.
Mr.\ Friedman brings the three things this round actually needs --- an
anchor customer, distribution, and capital-markets judgment --- because of
who he is and what he has already built.

\begin{itemize}[leftmargin=1.4em]
  \item \textbf{Track record.} Co-founder of \textbf{GDA Capital}, a
  blockchain merchant bank and venture studio with a portfolio spanning
  Web3 infrastructure, fintech, and the metaverse (Flashy Finance, Alpaca
  Network, Metaverse Group, NFT~BAZL, U-Topia, and others). Co-founder of
  \textbf{Secure Digital Markets}, a global institutional digital-asset
  brokerage through which he has facilitated \$2B+ in transactions with
  institutions, family offices, governments, exchanges, and Fortune~500
  counterparties. Founder of Metaverse Group and NFT~BAZL; earlier, a
  consumer-tech operator with an exit to WeWork.
  \item \textbf{Already aligned.} He invested personally into a Lux Network
  validator --- protocol-level, at-risk exposure to the network ahead of
  this round. Conviction already backed with capital is the strongest
  possible starting point for a lead.
  \item \textbf{Anchor customer.} SDM is the natural flagship white-label
  venue (\S\ref{sec:use} discussion). A lead who operates the first
  regulated deployment collapses the hardest B2B risk --- landing the first
  customer --- into an opening condition, and aligns the round's success
  with a business he already runs.
  \item \textbf{Distribution and capital.} GDA is a natural co-lead /
  syndicate source and brings the institutional, exchange, and family-office
  network that both the public launch (liquidity, listings, market-making)
  and the venue rollout (banking and fiat rails, counterparties) require.
  \item \textbf{The invitation.} Lead the \$100M round. A lead sets terms,
  anchors the allocation, and carries credibility into the syndicate; board
  or observer rights, information rights, and pro-rata are on the table.
\end{itemize}

%% ============================================================
\section{Risks and How They Are Managed}

Stated plainly, in the same spirit as the IP memorandum.

\begin{enumerate}[leftmargin=1.6em]
  \item \textbf{Execution / network operations.} Running a live public
  mainnet is operationally demanding; incidents happen and are recovered.
  The mitigant is a hardened validator set, post-quantum finality, and a
  disciplined recovery playbook --- but investors should underwrite an
  operating network, not a frictionless one.
  \item \textbf{Regulatory.} Broker-dealer, ATS, and transfer-agent
  licensing is slow and jurisdiction-specific; timelines are outside the
  company's sole control. Mitigant: jurisdiction-configurable architecture,
  a white-label model that lets \emph{licensed operators} carry the
  registration, and legal spend budgeted in the round.
  \item \textbf{IP / \S101 and prior art.} U.S.\ transaction-control claims
  carry abstract-idea risk; 2016--2018 state-channel literature is the
  invalidity hunting ground. Mitigants: the EPO granted the counterpart, the
  claims recite concrete architecture, and an invalidity-grade search
  precedes any acquisition price. The thesis is defensive depth + licensing,
  not assertion revenue.
  \item \textbf{Token / market.} Public launch outcomes depend on market
  conditions and liquidity depth. Mitigant: protocol-owned liquidity and
  market-making are funded line items, not afterthoughts.
  \item \textbf{Concentration.} A lead-anchored round concentrates early
  influence. Mitigant: standard governance, a broadened syndicate, and
  information rights for all round participants.
\end{enumerate}

%% ============================================================
\section{Next Steps}

\begin{enumerate}[leftmargin=1.6em]
  \item \textbf{Diligence pack.} Share the technical proof set (live mainnet
  endpoints, the on-chain DEX-settlement transaction, the running exchange
  stack) and the W3A IP memorandum with counsel.
  \item \textbf{Term conversation.} Agree instrument, valuation frame, lead
  allocation, and governance with Mr.\ Friedman.
  \item \textbf{IP closing path.} Run W3A formation and the '853 diligence
  (title, validity, family-completeness) in parallel, docketing the
  31~October~2026 maintenance window.
  \item \textbf{Syndicate.} With the lead set, open the balance of the round
  to aligned co-investors.
  \item \textbf{Launch sequencing.} Lock the public-launch and first
  regulated-venue timelines against round close.
\end{enumerate}

\vspace{1.5em}
\noindent\textit{This memorandum is confidential and prepared solely for
Zach Friedman's evaluation. It is not an offer to sell or a solicitation to
buy any security or token, and contains forward-looking statements that are
subject to risk and change. All financial figures are illustrative working
frames for discussion; valuation and terms are subject to definitive
documentation and diligence. Technical milestones described (live mainnet,
on-chain settlement, the running exchange stack) are current as of July
2026; the W3A patent acquisition is proposed and subject to the diligence
enumerated in its own memorandum.}

\end{document}
